\documentclass[UTF8,zihao=-4]{ctexart}
\usepackage[a4paper,margin=2.2cm]{geometry}
\usepackage{fontspec}
\usepackage{tabularx}
\usepackage{array}
\usepackage{ragged2e}
\usepackage{fancyhdr}
\usepackage{hyperref}
\setCJKmainfont{Noto Serif SC}
\setCJKsansfont{Noto Sans SC}
\setmainfont{Noto Serif SC}
\setlength{\parindent}{0pt}
\setlength{\parskip}{0.45em}
\pagestyle{fancy}
\fancyhf{}
\fancyhead[L]{商务智能}
\fancyfoot[C]{\thepage}
\hypersetup{colorlinks=true,linkcolor=black,urlcolor=blue}
\begin{document}
\title{6、操作型/分析型数据对比}
\author{}
\date{}
\maketitle

\begin{tabularx}{\linewidth}{|*{3}{>{\RaggedRight\arraybackslash}X|}}
\hline
特性 & 操作型数据（DB） & 分析型数据（DW） \\
\hline
定位 & 面向应用的事务处理 & 面向主题的数据分析 \\
\hline
DB设计 & E-R模型 & 星型/雪花模型，数据立方体 \\
\hline
数据 & 当前的、最新的 & 历史的，具有时间跨度 \\
\hline
汇总 & 原始的，细节的 & 集成的，一致的 \\
\hline
视图 & 详细的，关系的 & 总体的，多维的 \\
\hline
操作类型 & 读/写（可变的） & 读（稳定的） \\
\hline
存取请求 & 可预知的 & 事先未知的 \\
\hline
访问记录 & 一次操作少量记录 & 一次操作大量记录 \\
\hline
DB规模 & 100MB \textasciitilde{} GB & TB \\
\hline
工作单位 & 短的，简单事务 & 复杂查询 \\
\hline
性能要求 & 对性能要求高 & 对性能要求较宽松 \\
\hline
\end{tabularx}

\end{document}
